\section{Problem Formulation}

\textbf{Task setting.} Retrieval-augmented question answering. Given a question $q$, a
retriever returns documents $D=\{d_1,\dots,d_k\}$, and an LLM produces an answer $a$ that
cites document ids. We study the downstream \emph{verification} task: for each atomic
claim in $a$, decide whether it is supported by the specific source it cites. The base RAG
system (vector-DB builder $+$ retrieval-augmented QA with citations) is the
course-provided Program~1/Program~2, so our work targets the \emph{evaluation of published
verification techniques} on top of it, not re-implementing retrieval or generation.

\textbf{Information available.} The verifier sees the answer $a$, the retrieved passages
$D$ (full text $+$ ids), and the citation markers in $a$. It does not see a gold answer.

\textbf{Supervision / feedback signal.} Human span-level hallucination annotations in
RAGTruth~\citep{niu2024ragtruth} serve as ground truth; a small held-out, human-labeled
set of (claim, cited-passage, grounded?) pairs is used only to configure each published
scorer's decision threshold where applicable.

\textbf{Optimization / evaluation objective.} For each published scorer
$g(\text{claim},\text{source})$ under study, we measure how well its supported/unsupported
decision matches human labels (macro-F1), alongside its cost, latency, and run-to-run
determinism. The study's aim is comparative: identify which published technique(s) give
the best accuracy/cost/reproducibility trade-off as a RAG verification layer.

\textbf{Assumptions.} (i) Answers can be segmented into reasonably atomic claims by
published methods or simple heuristics. (ii) The compared scorers are used as published
(off-the-shelf), with only documented configuration. (iii) Modest compute: a single
consumer GPU and public/local models.
